\documentclass[12pt,a4paper]{article}
\usepackage[russian]{babel}
\usepackage{fontspec}
\usepackage{geometry}
\usepackage{booktabs}
\usepackage{longtable}
\usepackage{array}
\usepackage{multirow}
\usepackage{enumitem}
\usepackage{titlesec}
\usepackage{xcolor}
\usepackage{siunitx}

\usepackage{pdfpages}

\usepackage{tablefootnote}

% Дополнительные пакеты для данного документа
\usepackage{caption}
\usepackage{subcaption}
\usepackage{listings}
% \usepackage{hyperref}

\geometry{a4paper, margin=25mm}
\setmainfont{Liberation Serif}
\setsansfont{Liberation Sans}

\definecolor{adblue}{RGB}{0,84,166}

\titleformat{\section}{\large\bfseries\color{adblue}}{}{0em}{}[\titlerule]
\titleformat{\subsection}{\normalsize\bfseries}{}{0em}{}
\titleformat{\subsubsection}{\normalsize\itshape}{}{0em}{}

\sisetup{output-decimal-marker={,}}

\usepackage{tikz}
\usetikzlibrary{positioning, calc, backgrounds, math, 3d, perspective, arrows.meta, graphs, graphdrawing, fit, shapes.geometric, trees}
    

\usepackage{pgfplots} % для построения графиков
\pgfplotsset{compat=1.18} % версия для совместимости
\usepackage{tkz-euclide}

% --- Подключение пакета алгоритмов (без флага algo2e, чтобы вернуть \begin{algorithm}) ---
\usepackage[ruled,vlined]{algorithm2e}
%\usepackage{caption}
\newcommand{\Gray}[1]{\textcolor{gray}{#1}} 
% Настройка подписей
\DeclareCaptionLabelFormat{gost}{#1 #2}
\captionsetup[table]{labelformat=gost,labelsep=endash,justification=justified,singlelinecheck=false,font=normal}
\captionsetup[figure]{labelformat=gost,labelsep=endash,justification=centering,font=normal}

% Локализация algorithm2e на русский (ОБЯЗАТЕЛЬНО после загрузки пакета)
\SetAlgorithmName{Алгоритм}{}{}
\SetAlgoCaptionSeparator{---}
\SetKwInput{Input}{Вход}
\SetKwInput{Output}{Выход}

\usepackage{url}
% Говорим пакету использовать моноширинный шрифт (как у texttt)
\Urlmuskip=0mu plus 1mu % гибкие пробелы для формул, если нужно
\DeclareUrlCommand\code{\urlstyle{tt}} 

% Библиотека схемотехники
\usepackage[siunitx, RPvoltages, american]{circuitikzgit}
\ctikzset{
    amplifiers/fill=cyan,
    sources/fill=green!20,
    diodes/fill=white,
    resistors/fill=violet,
    grounds/scale=1.2,
}

\usepackage{iftex}

\ifXeTeX % Если используется TeTeX или TeLaTeX
  \typeout{Режим XeTeX}
  \usepackage{xltxtra}
  \usepackage{fontspec}
  \usepackage{polyglossia}
  \setmainlanguage{russian}
  \setotherlanguage{english}
  \setkeys{russian}{babelshorthands=true} % По идее включает <<, >>, -- ---

  \setmainfont{Times New Roman}[Ligatures=TeX]
  \setromanfont{Times New Roman}[Ligatures=TeX]
  \setsansfont{Arial}[Ligatures=TeX]
  \setmonofont{Courier New}[Ligatures=TeX] 

  \newfontfamily{\cyrillicfont}{Times New Roman}
  \newfontfamily{\cyrillicfontrm}{Times New Roman}
  \newfontfamily{\cyrillicfonttt}{Courier New}
  \newfontfamily{\cyrillicfontsf}{Arial}
\fi

\ifLuaTeX % Если используется LuaLatex
  \typeout{Режим LuaTeX}
%  \usepackage{luatextra}
  \usepackage{fontspec}
  \defaultfontfeatures{Ligatures={TeX}}
  \usepackage{polyglossia}
  \setmainlanguage{russian}
  \setotherlanguage{english}
  \setmainfont{Times New Roman}[Script=Cyrillic, Ligatures=TeX]
  \setromanfont{Times New Roman}[Script=Cyrillic, Ligatures=TeX]
  \setsansfont{Arial}[Script=Cyrillic, Ligatures=TeX]
  \setmonofont{Courier New}[Ligatures=TeX]
  \newfontfamily{\cyrillicfont}{Times New Roman}[Script=Cyrillic, Ligatures=TeX]
  \newfontfamily{\cyrillicfontrm}{Times New Roman}[Script=Cyrillic, Ligatures=TeX]
  \newfontfamily{\cyrillicfonttt}{Courier New}
  \newfontfamily{\cyrillicfontsf}{Arial}[Script=Cyrillic, Ligatures=TeX]

  \usegdlibrary{trees}
  \usegdlibrary{force}
\fi

\ifPDFTeX % Если используется 8битный PDFTex
  \typeout{8-битный режим pdfTex}
  % For pdfTeX we must use old
  % 8-bit font technologies
  \usepackage[T2A]{fontenc}
  \usepackage[english,russian]{babel}
  \usepackage[utf8]{inputenc}
  \usepackage[english, russian]{babel}

  %\usepackage{newtxmath}
  \usepackage{amsmath}
  \usepackage{amsthm}
  \usepackage{amssymb}
  \usepackage{amsfonts}
  \usepackage{mathtext}
\fi

\usepackage{amsthm}  % Возможности AMSMath
\usepackage{amsmath}
\usepackage{amssymb}

\usepackage{esvect} % Для красивых стрелок
\usepackage{bm} % Для жирных символов

%Hyphenation rules
\usepackage{hyphenat}
\hyphenation{ма-те-ма-ти-ка вос-ста-нав-ли-вать ото-бра-жа-ют-ся сге-не-ри-ро-ван-ном сим-во-лы}
% \usepackage[english, russian]{babel}

\usepackage{graphicx}
\graphicspath{ {./images/} }

\usepackage{empheq}

\usepackage{float}

% \usepackage{comment}

\usepackage[hidelinks,
    unicode=true,          % поддерживает Unicode в закладках
    psdextra=true,      % автоматически конвертирует простую математику в текст
    % focusdefault=false
]{hyperref}


% Окружение "Теорема"
\newtheorem{theorem}{Теорема}[section]
\newtheorem{corollary}{Corollary}[theorem]
\newtheorem{lemma}[theorem]{Lemma}

% Окружение "Определение"
\theoremstyle{definition} % Прямой шрифт, нумерация
\newtheorem{definition}{Определение}[section]
\newtheorem{example}{Пример}[section]

% Окружение "Замечание"
\theoremstyle{remark} % Прямой шрифт, без нумерации
% Окуржение "Решение"
\newtheorem*{solution}{Решение}
\newtheorem*{remark}{Замечание}

\usepackage{polynom}
\usepackage{tabularx}
\usepackage{cancel}

\addto\captionsrussian{%
  \renewcommand{\figurename}{Рис.}%
  \renewcommand{\tablename}{Табл.}%
}

\renewcommand{\arraystretch}{1.5} % Высота строки таблицы

\renewcommand{\labelitemi}{\(\circ\)} % Бюллет -- пустой кружок в itemize


\graphicspath{{figures/} {images/}}

\title{ГЛАВА 3. ВЫЧИСЛИТЕЛЬНЫЙ ЭКСПЕРИМЕНТ ПО СОВМЕСТНОМУ ВОССТАНОВЛЕНИЮ ПРОСТРАНСТВЕННО-МОЗАИЧНЫХ RAW-ИЗОБРАЖЕНИЙ И СВЕРХРАЗРЕШЕНИЮ}
\author{Наталия А. Рыбкина}
\date{\today}

\begin{document}

\maketitle

\tableofcontents

\section{ГЛАВА 3. ВЫЧИСЛИТЕЛЬНЫЙ ЭКСПЕРИМЕНТ ПО СОВМЕСТНОМУ ВОССТАНОВЛЕНИЮ ПРОСТРАНСТВЕННО-МОЗАИЧНЫХ RAW-ИЗОБРАЖЕНИЙ И СВЕРХРАЗРЕШЕНИЮ}

\subsection{3.1. Методология моделирования сквозного оптико-цифрового тракта деградации входных данных}

Для прецизионного обучения и верификации разработанной нейросетевой архитектуры совместной интерполяции и сверхразрешения (\textit{Joint Demosaicing and Super-Resolution}, $\operatorname{JDSR}$) сформирована физически адекватная модель деградации высококачественных RGB-изображений высокой чёткости ($Y \in \mathbb{R}^{H \times W \times 3}$) в низкоразмерные пространственно-мозаичные RAW-пакеты данных ($X \in \mathbb{R}^{\frac{H}{4} \times \frac{W}{4} \times 4}$). 

Математически процесс формирования входного сигнала сенсора в рамках разработанного программного комплекса описывается следующим сквозным операторным уравнением:
\begin{equation}
X = \mathcal{E}_{\operatorname{subchannels}} \left( \mathcal{M}_{\operatorname{Bayer}} \left( \mathcal{N}_{\operatorname{matrix}} \left( \mathcal{H}_{\operatorname{PSF}} \left( \mathcal{D}_{\downarrow \gamma} (Y) \right) \right) \right) \right)
\end{equation}
где $\mathcal{D}_{\downarrow \gamma}$ --- оператор оптического ресэмплинга с масштабным коэффициентом $\gamma$, $\mathcal{H}_{\operatorname{PSF}}$ --- функция рассеяния точки оптической системы объектива, $\mathcal{N}_{\operatorname{matrix}}$ --- оператор аддитивно-мультипликативной стохастической деградации, $\mathcal{M}_{\operatorname{Bayer}}$ --- пространственно-вариантная маска мозаики цветных фильтров, а $\mathcal{E}_{\operatorname{subchannels}}$ --- оператор пространственной декомпозиции фазовых подканалов КМОП-матрицы ($\text{CMOS}$).

\subsubsection{3.1.1. Математическое описание оператора пространственного и фазового маскирования}

В процессе генерации низкокачественных входных данных ($\text{LQ}$) из непрерывного RGB-потока высокого разрешения ($\text{HQ}$) извлекается центральная пространственная область фиксированной геометрии $W_{\operatorname{HQ}} \times H_{\operatorname{HQ}}$ (в текущей конфигурации вычислительного стенда --- $512 \times 512$ пикселей). Для полного исключения деструктивных фазовых сдвигов цветовой поднесущей мозаики Байера при последующих иерархических дискретных преобразованиях, пространственные координаты левого верхнего угла кадрирования $(x_0, y_0)$ принудительно квантуются с шагом, кратным базовому периоду суперпикселя:
\begin{equation}
x_0 \equiv 0 \pmod 4, \quad y_0 \equiv 0 \pmod 4
\end{equation}
Данное математическое ограничение гарантирует инвариантность топологической структуры цветовой сетки на всех этапах цифровой деградации и исключает геометрический рассинхронизм парных патчей при обучении.

\subsubsection{3.1.2. Аппроксимация функции рассеяния точки оптической системы (PSF)}

Имитация ограниченной частотно-контрастной характеристики ($\text{ЧКХ}$) реального объектива и дифракционного предела регистрирующей системы реализуется в два последовательных этапа:
\begin{enumerate}
    \item \textbf{Геометрический ресэмплинг ($\mathcal{D}_{\downarrow \gamma}$):} Исходная матрица масштабируется с коэффициентом $\gamma = 2$ (параметр конфигурации \texttt{process\_data.downscale\_factor}) посредством билинейного оператора дискретизации, уничтожая первичную субпиксельную информацию:
    \begin{equation}
    I_{\text{small}} = \mathcal{D}_{\downarrow 2}(Y), \quad I_{\text{small}} \in \mathbb{R}^{\frac{H}{2} \times \frac{W}{2} \times 3}
    \end{equation}
    \item \textbf{Размытие пятном рассеяния ($\mathcal{H}_{\operatorname{PSF}}$):} Оптические аберрации моделируются операцией двумерной дискретной свёртки промежуточного распределения освещённости с изоморфным гауссовым ядром ($\text{Function Scatter Point}$, $\text{PSF}$):
    \begin{equation}
    I_{\text{blur}}(x,y) = I_{\text{small}}(x,y) \ast G(x,y, \sigma), \quad G(x,y,\sigma) = \frac{1}{2\pi\sigma^2}\exp\left(-\frac{x^2+y^2}{2\sigma^2}\right)
    \end{equation}
    где параметр среднеквадратического отклонения ядра зафиксирован на уровне $\sigma = 0{,}5$ (параметр конфигурации \texttt{psf\_sigma}) для симуляции реального пятна рассеяния объектива.
\end{enumerate}

\subsubsection{3.1.3. Стохастическое моделирование стохастического и пространственно-коррелированного шума КМОП-матрицы}

Для обеспечения адекватности синтезированных данных реальным физическим процессам регистрации светового потока полупроводниковыми структурами, в операторе $\mathcal{N}_{\operatorname{matrix}}$ реализовано двухкомпонентное моделирование шумовой компоненты $\eta$ с жестко заданным соотношением сигнал/шум ($\text{SNR} = 20~дБ$, параметр конфигурации \texttt{snr\_db}):
\begin{enumerate}
    \item \textbf{Некоррелированный белый гауссов шум:} Моделирует тепловой шум считывания КМОП-сенсора и дробовой фотонный шум с независимым распределением в каждом пространственном элементе матрицы ($\eta_{\text{white}} \sim \mathcal{N}(0, \sigma_n^2)$).
    \item \textbf{Пространственно-коррелированный шум:} Имитирует прохождение стохастических помех сквозь сглаживающие оптические фильтры ($\text{Anti-Aliasing}$) цифровых камер. Формируется путём математического пропускания первичного белого шума через фильтр свёртки с гауссовым ядром: $\eta_{\text{corr}} = \eta_{\text{white}} \ast G(x,y,\sigma)$.
\end{enumerate}

\subsubsection{3.1.4. Топологические инварианты структуры Bayer-RGGB при операциях пространственной декомпозиции субканалов}

Перевод непрерывного распределения энергии по цветовым каналам в дискретный формат RAW-данных осуществляется оператором проекции на мозаичную решётку цветных фильтров $\mathcal{M}_{\operatorname{Bayer}}$ по каноническому периодическому шаблону $\text{RGGB}$:
\begin{equation}
B_{\text{float}}(x,y) = \sum_{c \in \{R,G,B\}} I_{\text{blur}}(x,y,c) \cdot \delta\left(c - \mathcal{P}_{\operatorname{RGGB}}(x \pmod 2, y \pmod 2)\right)
\end{equation}
где $\mathcal{P}_{\operatorname{RGGB}}$ --- матрица фаз светофильтров, а $B_{\text{float}} \in \mathbb{R}^{\frac{H}{2} \times \frac{W}{2} \times 1}$ --- одноканальный мозаичный массив плавающей точки в диапазоне $[0,\, 1]$.

Для обеспечения пространственной стационарности свёрточных слоёв глубокой сети, в модуле $\mathcal{E}_{\operatorname{subchannels}}$ реализован алгоритм декомпозиции мозаики в четырёхканальный тензор фазовых плоскостей $X \in \mathbb{R}^{\frac{H}{4} \times \frac{W}{4} \times 4}$. Распределение пикселей для каждой локальной координаты $(x, y)$ подчиняется строгой топологической привязке:
\begin{equation}
\begin{cases}
X(x, y, 0) = B_{\text{float}}(2x, 2y) \quad &\text{(Канал R --- чётные строки, чётные столбцы)}, \\
X(x, y, 1) = B_{\text{float}}(2x, 2y+1) \quad &\text{(Канал $G_1$ --- чётные строки, нечётные столбцы)}, \\
X(x, y, 2) = B_{\text{float}}(2x+1, 2y) \quad &\text{(Канал $G_2$ --- нечётные строки, чётные столбцы)}, \\
X(x, y, 3) = B_{\text{float}}(2x+1, 2y+1) \quad &\text{(Канал B --- нечётные строки, нечётные столбцы)}.
\end{cases}
\end{equation}
Данный шаг декомпозиции приводит к закономерному уменьшению геометрических размеров каждой хроматической плоскости в 2 раза, формируя сквозной коэффициент пространственной деградации $F_{\operatorname{total}} = \gamma \cdot 2 = 4$. Соответственно, для компенсации деградации и релаксации обратной некорректной задачи реставрации, на входном слое нейросетевой обёртки зафиксирован оператор субпиксельного сдвига признаков \texttt{PixelShuffle(4)} ($\texttt{network\_g.upscale\_factor} = 4$).

\subsection{3.2. Архитектурная адаптация нелинейно-безактивационной нейросети NAFNet для задач JDSR}

Для решения совмещённой задачи демозаикации и четырёхкратного сверхразрешения ($\operatorname{JDSR}$) стандартные свёрточные архитектуры, оперирующие в пространстве яркостных сигналов высокого уровня, оказываются неэффективными. Это обусловлено тем, что классические нелинейные функции активации (такие как $\operatorname{ReLU}$ или $\operatorname{GELU}$) обладают свойством жесткого порогового отсечения отрицательных флуктуаций значений. В условиях работы с низкоуровневыми сигналами КМОП-сенсора ($\text{RAW}$-данными) подобные ограничения приводят к искусственному занулению стохастической шумовой компоненты и необратимому разрушению фазовых соотношений между каналами мозаики Байера. 

В рамках данного исследования в качестве базового вычислительного ядра конвейера реставрации принята и модифицирована концепция нелинейно-безактивационной нейросети $\operatorname{NAFNet}$ (\textit{Nonlinear Activation Free Network}) \cite{chen2022simple}, адаптированная под сквозную топологию пространственно-фазовых преобразований.

\subsubsection{3.2.1. Концепция Simple Gate и сохранение физической линейности RAW-сигнала}

Фундаментальной теоретической особенностью архитектуры $\operatorname{NAFNet}$ является полный отказ от классических нелинейных функций активации в пользу операции контролируемого линейного масштабирования признаков, реализуемой через блок $\operatorname{Simple~Gate}$ ($\operatorname{SG}$). Математически данное преобразование описывается как поэлементное произведение двух разностных проекций тензора признаков, полученных путём разделения канальной размерности:
\begin{equation}
\operatorname{SG}(X_1, X_2) = X_1 \odot X_2
\end{equation}
где $\odot$ --- оператор произведения Адамара, а $X_1, X_2 \in \mathbb{R}^{H \times W \times \frac{C}{2}}$ --- комплементарные матрицы латентного пространства признаков, сформированные после прохождения слоя групповой нормализации ($\operatorname{LayerNorm}$). 

Поскольку оператор $\operatorname{SG}$ является непрерывным, гладким и дифференцируемым на всей числовой оси, он эффективно аппроксимирует нелинейные межканальные взаимодействия, сохраняя при этом строгую линейность исходного физического сигнала сенсора. Отсутствие сингулярностей в нуле исключает появление эффекта «застревания» градиентов на ранних этапах обучения, что позволяет глубоким свёрточным блокам сети оптимизировать весовые коэффициенты непосредственно в терминах фазовых сдвигов и пространственно-периодических плотностей спектра Фурье.

\subsubsection{3.2.2. Синтез пространственно-фазовых входных адаптеров на основе субпиксельного сдвига (PixelShuffle)}

Оригинальная архитектура сети $\operatorname{NAFNet}$ спроектирована для реставрации трёхканальных $\operatorname{RGB}$-изображений без изменения их геометрических размеров ($1\times$). Для реализации сквозного конвейера четырёхкратного апскейла ($4\times$) и интерполяции мозаики без внесения промежуточных артефактов («замыливания» линейными фильтрами) в программном комплексе разработана специализированная топологическая обёртка $\texttt{NAFNetDemosaicSuperResolutionWrapper}$. 

Схема циркуляции потоков данных внутри модифицированного графа вычислений формализуется следующим образом:
\begin{enumerate}
    \item \textbf{Входной пространственно-фазовый адаптер:} На вход подаётся сформированный на этапе деградации четырёхканальный тензор низкого разрешения $X_{\text{LQ}} \in \mathbb{R}^{\frac{H}{4} \times \frac{W}{4} \times 4}$, инкапсулирующий изолированные фазы $\text{RGGB}$. Первичный свёрточный блок $\operatorname{Conv2d}$ транслирует его в латентное подпространство с размерностью каналов $C_{\text{mid}} = 48$.
    \item \textbf{Оператор субпиксельного сдвига:} Пространственная структура тензора мгновенно перестраивается с помощью математического оператора $\operatorname{PixelShuffle}$ с масштабным коэффициентом $\mathcal{P}_{\operatorname{Shuffle}}(4)$. Это позволяет перепроецировать скрытые фазовые признаки из низкоразмерного пространства субканалов в полноразмерную координатную сетку целевого $\operatorname{RGB}$-изображения высокого разрешения ($H \times W \times 3$).
    \subsubsection*{Методическая схема прохождения тензора:}
    \begin{equation}
    X_{\text{LQ}} \left(\frac{H}{4} \times \frac{W}{4} \times 4\right) \xrightarrow{\operatorname{Conv2d}} X_{\text{feat}} \left(\frac{H}{4} \times \frac{W}{4} \times 48\right) \xrightarrow{\mathcal{P}_{\operatorname{Shuffle}}(4)} X_{\text{HQ\_space}} (H \times W \times 3)
    \end{equation}
    \item \textbf{Магистральный экстрактор признаков:} Глубокая сеть $\operatorname{NAFNet}$ (состоящая из каскада блоков $\operatorname{Simple~Gate}$ и локального межканального внимания) функционирует непосредственно на полной геометрической сетке $H \times W$, оперируя внутренним числом каналов $C_{\text{int}} = 4$. Таким образом, скрытые слои осуществляют сквозной нелинейный синтез утерянных цветовых компонент, используя межканальные корреляции на макро- и микроструктурных уровнях.
    \item \textbf{Выходной адаптер:} Финальный свёрточный слой выполняет проекцию из четырёхмерного пространства внутренних признаков сети в ортогональный трёхмерный базис цветового пространства $\operatorname{RGB}$ высокого качества ($\hat{Y} \in \mathbb{R}^{H \times W \times 3}$).
\end{enumerate}

\subsubsection{3.2.3. Алгоритм селективного частичного переноса параметров (Partial Weight Transfer) из латентных пространств SIDD}

Для преодоления проблемы «холодного старта» вычислительного графа в условиях ограниченного объёма видеопамяти графического ускорителя $\text{NVIDIA GeForce RTX 4070 Ti}$ ($12~\text{ГБ}~\text{VRAM}$) разработан и внедрён алгоритм селективной инициализации параметров на основе частичного переноса весов (\textit{Partial Weight Transfer}). 

В качестве донорной модели используется глубокая сеть $\operatorname{NAFNet}$, предварительно обученная на крупномасштабном датасете $\text{SIDD}$ (\textit{Smartphone Image Denoising Dataset}) для задачи фильтрации высокочастотных шумов. Из-за топологической перестройки входного и выходного интерфейсов разработанной обёртки, свёрточные слои первичного приёма данных ($\texttt{intro}$) и финальной проекции ($\texttt{ending}$) имеют несовместимую геометрию тензоров весов и принудительно исключаются из процедуры инициализации. Исполняемый скрипт осуществляет автоматический послойный анализ структуры и переносит только совместимые конфигурации:
\begin{equation}
\Theta_{\text{target}} = \left\{ \theta_i \in \Theta_{\text{SIDD}} \;\middle|\; \operatorname{shape}(\theta_i)_{\text{target}} == \operatorname{shape}(\theta_i)_{\text{SIDD}} \right\}
\end{equation}
Данная процедура позволила успешно импортировать весовые коэффициенты $661$ внутреннего слоя экстракции признаков. Экспериментально доказано, что использование предобученного латентного пространства фильтрации шумов в качестве начального приближения существенно снижает общую энтропию системы на старте обучения. Это обеспечивает стабильную и высокоскоростную сходимость сложной многокритериальной функции потерь без возникновения градиентных всплесков и шоковых осцилляций весов.


\subsection{3.3. Разработка многокритериального функционала потерь с динамическим частотным прогревом графа}

Для обеспечения прецизионной сходимости разработанного нейросетевого конвейера и одновременного контроля как низкочастотных полутоновых переходов (общая геометрия сцены, макроконтуры, баланс освещённости), так и высокочастотных пространственно-периодических структур (микротекстура лепестков люпинов, ворсинки на стеблях растений) синтезирована многокритериальная функция потерь $\mathcal{L}_{\operatorname{total}}$. Временной контур оптимизации параметров модели разделён на две фазы с помощью разрывного оператора Хевисайда, что позволяет полностью купировать эффекты градиентного шока на начальных этапах эволюции весовых коэффициентов.

\subsubsection{3.3.1. Пространственная компонента оптимизации на основе L1-нормы}

В качестве базисного пространственного критерия в программном комплексе используется средняя абсолютная ошибка ($L_1$-norm, \textit{Mean Absolute Error}, $\operatorname{MAE}$), вычисляемая между матрицей предсказанного моделью высокого разрешения $\operatorname{RGB}$-изображения $\hat{Y} = f(X_{\text{LQ}}; \Theta)$ и эталонным кадром оригинального качества $Y$:
\begin{equation}
\mathcal{L}_{L1}(\Theta) = \frac{1}{H \cdot W \cdot C} \sum_{x=1}^{H} \sum_{y=1}^{W} \sum_{c=1}^{C} \left| \hat{Y}(x,y,c) - Y(x,y,c) \right|
\end{equation}
Данный функционал минимизирует попиксельную разность амплитуд в пространственной области $\mathbb{R}^2$. Обладая высокой устойчивостью к стохастическим выбросам и аномалиям выборки, критерий $\mathcal{L}_{L1}$ обеспечивает стабильное выравнивание средних яркостных характеристик кадра. Однако из-за математического свойства усреднения градиентов по площади, изолированное применение пространственной нормы приводит к эффекту локального сглаживания («замыливания») межкомпонентных границ на экстремальных частотах спектра.

\subsubsection{3.3.2. Логарифмический фокально-частотный критерий (FFL-Log) в частотной области Фурье}

Для преодоления ограничений пространственной фильтрации, компенсации спектрального наложения (\textit{aliasing}) и ультимативного подавления цветового муара Байера, в контур автодифференцирования введён нелинейный фокально-частотный функционал потерь \cite{jiang2021focal}, модифицированный автором в логарифмическую шкалу. 

Процедура вычисления спектрального рассогласования реализуется путём трансляции предсказанного и эталонного тензоров в комплексную частотную область посредством двумерного дискретного быстрое преобразования Фурье ($\text{2D-DFT}$):
\begin{equation}
\mathcal{F}(\hat{Y})_{u,v} = \sum_{x=0}^{M-1} \sum_{y=0}^{N-1} \hat{Y}_{x,y} \cdot \exp\left( -j 2\pi \left( \frac{ux}{M} + \frac{vy}{N} \right) \right)
\end{equation}
Для централизации постоянной составляющей яркости к геометрическому центру комплексной плоскости применяется оператор циклического сдвига четвертей спектра $\operatorname{Shift}(\cdot)$. С целью компрессии колоссального динамического диапазона амплитуд и частотного выравнивания градиентных потоков вводится нелинейное логарифмическое сжатие с константой масштабирования $\gamma = 100.0$:
\begin{equation}
\mathcal{A}_{\log}(\hat{Y}) = \frac{\ln\left(1 + \gamma \cdot \left|\operatorname{Shift}\left(\mathcal{F}(\hat{Y})\right)\right|\right)}{\ln(1 + \gamma)}
\end{equation}
Для стабилизации численных вычислений и предотвращения деления на нуль при вычислении дифференциалов, константа логарифмического знаменателя принудительно зафиксирована в аппаратном буфере вычислительного графа через механизм \texttt{register\_buffer}.

Итоговый критерий $\mathcal{L}_{\operatorname{FFL-Log}}$ взвешивается пространственно-вариантной матрицей фокальных весов $w(u,v)$, которая динамически увеличивает штраф за ошибки на тех частотах, где рассогласование между моделью и эталоном максимально:
\begin{equation}
\mathcal{L}_{\operatorname{FFL-Log}}(\Theta) = \frac{1}{U \cdot V} \sum_{u=1}^{U} \sum_{v=1}^{V} w(u,v) \cdot \left( \mathcal{A}_{\log}(\hat{Y})_{u,v} - \mathcal{A}_{\log}(Y)_{u,v} \right)^2
\end{equation}
\begin{equation}
w(u,v) = \left| \left|\operatorname{Shift}\left(\mathcal{F}(\hat{Y})\right)\right|_{u,v} - \left|\operatorname{Shift}\left(\mathcal{F}(Y)\right)\right| Holy_{u,v} \right|^\alpha, \quad \alpha = 1
\end{equation}

\subsubsection{3.3.3. Динамическое управление контуром градиентов через разрывную ступенчатую функцию Хевисайда}

На начальных итерациях обучения, когда входной адаптер и весовые коэффициенты слоя субпиксельного сдвига $\mathcal{P}_{\operatorname{Shuffle}}(4)$ находятся в состоянии хаотической квазислучайной инициализации, спектральные градиенты Фурье генерируют высокоамплитудный шум. Прямая подача этих потоков в обратный проход приводит к сингулярностям весов и мгновенному развалу латентного пространства признаков.

Для обеспечения вычислительной устойчивости системы в диссертационном исследовании разработана методология динамического частотного прогрева (\textit{Dynamic Frequency Warm-up}). Полная целевая функция потерь $\mathcal{L}_{\operatorname{total}}$ параметризована текущим индексом эпохи обучения $t$ и формализуется с использованием разрывного дискретного оператора Хевисайда $\theta(x)$:
\begin{equation}
\mathcal{L}_{\operatorname{total}}(t) = \lambda_{1} \mathcal{L}_{L1} + \theta(t - t_{\operatorname{start}}) \cdot \lambda_{2} \mathcal{L}_{\operatorname{FFL-Log}}(t)
\end{equation}
\begin{equation}
\theta(t - t_{\operatorname{start}}) = \begin{cases} 
0, & \quad \text{при } t < t_{\operatorname{start}}, \\ 
1, & \quad \text{при } t \ge t_{\operatorname{start}}. 
\end{cases}
\end{equation}
Эмпирически верифицированные весовые множители зафиксированы в пропорции $\lambda_1 = 1{,}5$ и $\lambda_2 = 0{,}2$. Временной порог переключения контуров оптимизации жёстко детерминирован в YAML-конфигурации на уровне $t_{\operatorname{start}} = \texttt{losses.ffl\_start\_epoch} = 20$. 

Данная стратегия разделения частотно-временных диапазонов обеспечивает строго двухэтапную эволюцию параметров модели: первые 20 эпох граф вычислений стабильно прогревается на пространственной норме $\mathcal{L}_{L1}$, настраивая базовую геометрию и фазы `PixelShuffle`, а начиная с 21-й эпохи функция Хевисайда автоматически подключает логарифмический спектральный лосс, который осуществляет финальный бескомпромиссный дожим байеровского муара и реставрацию субпиксельных деталей.

\subsection{3.4. Программная оркестрация, метрический аппарат и стендовые условия верификации}

Для проведения численных симуляций, обучения разработанных архитектурных модификаций и статистического анализа результатов реставрации в рамках диссертационного исследования был спроектирован и развёрнут специализированный программно-аппаратный комплекс. Оркестрация вычислительных потоков, распределение ресурсов видеопамяти и верификация фазовой когерентности полностью автоматизированы на уровне разработанных модулей среды программирования $\text{Python 3.10}$ с использованием фреймворка глубокого обучения $\text{PyTorch 2.1}$.

\subsubsection{3.4.1. Спецификация вычислительного стенда (Ubuntu под WSL2, NVIDIA GeForce RTX 4070 Ti 12GB VRAM)}

Экспериментальный мейл-ран и сопутствующие нагрузочные тесты выполнялись на базе гетерогенной вычислительной инфраструктуры со следующими техническими характеристиками:
\begin{itemize}
    \item \textbf{Центральный процессор и оперативная память:} хост-система под управлением архитектуры $\text{x86-64}$, оснащённая $32~\text{ГБ}$ системной памяти $\text{RAM}$ функционирующей в многоканальном режиме.
    \item \textbf{Программная прослойка виртуализации ядра:} операционная система $\text{Ubuntu 22.04 LTS}$, развёрнутая внутри подсистемы $\text{WSL2}$ (\textit{Windows Subsystem for Linux}) на микроядре $\text{Windows 11}$. Данная конфигурация обеспечивает прямой бесконфликтный проброс системных вызовов к драйверам графического ускорителя через программные интерфейсы $\text{CUDA Direct}$.
    \item \textbf{Графический ускоритель:} дискретная видеокарта $\text{NVIDIA GeForce RTX 4070 Ti}$ на базе микроархитектуры $\text{Ada Lovelace}$, обладающая $12~\text{ГБ}$ аппаратной видеопамяти $\text{VRAM}$ и пропускной способностью шины $192~\text{бит}$. 
\end{itemize}

Для предотвращения критических ошибок нехватки динамической памяти (\textit{Out-Of-Memory, OOM}) в условиях жёсткого лимита в $12~\text{ГБ}$ $\text{VRAM}$, граф вычислений обёртки $\texttt{NAFNetDemosaicSuperResolutionWrapper}$ подвергся инженерной оптимизации. В частности, базовый функционал расчёта среднеквадратичной ошибки \texttt{nn.MSELoss()} был вынесен из итерационного цикла на уровень глобального модуля для статической аллокации памяти, а промежуточные операции копирования массивов через метод \texttt{.copy()} в NumPy-срезах были полностью исключены из потока валидации, что снизило пиковое потребление памяти на $\approx 18\%$.

\subsubsection{3.4.2. Алгоритмическое обеспечение модуля CustomNEFPairDataset и лимитирование пространственных аугментаций}

Синхронная загрузка и пространственно-координационная привязка парных изображений низкого ($\text{LQ}$) и высокого ($\text{HQ}$) разрешений реализуется разработанным классом датасета $\texttt{CustomNEFPairDataset}$. Для ликвидации геометрических сдвигов сцены и рассинхронизма контуров, в модуле осуществлён переход к динамическому масштабированию привязки координат срезов на основе единственного источника истины --- параметра $\texttt{network\_g.upscale\_factor} = 4$, считываемого напрямую из конфигурационного YAML-файла. Координаты случайного пространственного кропа низкого разрешения $(t_l, l_l)$ однозначно детерминируют положение патча высокого разрешения $(t_h, l_h)$ по закону:
\begin{equation}
t_h = t_l \cdot \texttt{upscale\_factor}, \quad l_h = l_l \cdot \texttt{upscale\_factor}
\end{equation}
Данный подход обеспечивает инвариантность загрузчика к любому пространственному масштабу апскейла и гарантирует, что границы кропа $\texttt{gt\_size} = 512$ строго соответствуют границам $\texttt{lq\_size} = 128$.

В процессе стохастического расширения обучающей выборки (\textit{Data Augmentation}) наложены жёсткие топологические ограничения на геометрические трансформации тензоров. В коде загрузчика принудительно заблокированы и исключены любые ортогональные повороты матриц на углы $90^\circ$ и $270^\circ$. Математическим обоснованием данного запрета является пространственная анизо anisotropy периодической решётки цветных фильтров Байера. Поворот координатной сетки на некратные $180^\circ$ углы инвертирует фазы каналов (например, фаза $R$ переходит в $G_1$ или $B$), нарушая топологическую привязку весов PixelShuffle и приводя к неустранимому искажению цветопередачи кадра. Допускается использование исключительно операторов зеркального отражения относительно горизонтальной ($\operatorname{hflip}$) и вертикальной ($\operatorname{vflip}$) осей симметрии, сохраняющих циклическую RGGB-структуру периодической матрицы.

\subsubsection{3.4.3. Аппарат двухкритериальной валидации (PSNR, SSIM) и метрики контроля стабильности (TV-ratio, Grad-variance)}

Оценка качества реставрации изображений выполняется в рамках изолированного потока класса $\texttt{TrainingValidator}$ по завершении каждой эпохи обучения. Для исключения краевых дефектов расчёт метрик производится строго на центрированных кропах (\textit{Center Crop}) с приведением непрерывных значений тензоров к дискретному 8-битному пространству целых чисел $\mathbb{Z} \in [0,\, 255]$.

Инструментальный контроль включает две базовые метрики качества:
\begin{enumerate}
    \item \textbf{Пиковое отношение сигнала к шуму ($\operatorname{PSNR}$):} фиксирует среднеквадратическую точность восстановления физического уровня сигнала яркости:
    \begin{equation}
    \operatorname{MSE} = \frac{1}{H \cdot W \cdot C}\sum_{x,y,c} \left( \hat{Y}_{\text{uint8}}(x,y,c) - Y_{\text{uint8}}(x,y,c) \right)^2, \quad \operatorname{PSNR} = 20 \cdot \log_{10} \left( \frac{255}{\sqrt{\operatorname{MSE}}} \right)
    \end{equation}
    \item \textbf{Индекс структурного сходства ($\operatorname{SSIM}$):} определяет степень сохранности геометрии, локального контраста и контуров объектов, коррелируя с экспертной оценкой восприятия человека. Расчёт выполняется поканально через интерфейсы \texttt{skimage.metrics} со сдвигом по третьей оси признаков (\texttt{channel\_axis=2}).
\end{enumerate}

Параллельно для верификации устойчивости градиентных потоков вычисляются две дополнительные метрики стабильности графа:
\begin{itemize}
    \item \textbf{Коэффициент полной вариации ($\operatorname{TV}_{\operatorname{ratio}}$):} нормированное отношение интегральных пространственных вариаций предсказанного и эталонного тензоров, контролирующее сохранение микроструктурного контраста и отсутствие эффектов сглаживания:
    \begin{equation}
    \operatorname{TV}_{\operatorname{ratio}} = \frac{\sum_{i} \sum_{j} |\hat{Y}_{i+1,j} - \hat{Y}_{i,j}| + \sum_{i} \sum_{j} |\hat{Y}_{i,j+1} - \hat{Y}_{i,j}|}{\left(\sum_{i} \sum_{j} |Y_{i+1,j} - Y_{i,j}| + \sum_{i} \sum_{j} |Y_{i,j+1} - Y_{i,j}|\right) + \epsilon}
    \end{equation}
    \item \textbf{Дисперсия норм градиентов параметров ($\operatorname{Grad}_{\operatorname{variance}}$):} математический индикатор контролируемости сходимости, отслеживающий отсутствие затухания или взрыва градиентного шага по всей иерархии слоёв:
    \begin{equation}
    \operatorname{Grad}_{\operatorname{variance}} = \operatorname{Var}\left( \left\{ \|\nabla_{\theta_k} \mathcal{L}_{\operatorname{total}}\|_2 \right\}_{k=1}^{K} \right)
    \end{equation}
\end{itemize}
Синхронизированный валидационный лог через класс $\texttt{TrainingLogger}$ осуществляет запись метрик в TensorBoard (\texttt{self.tb\_logger}), а также транслирует сквозную аналитику в структурированный текстовый CSV-файл со строгим ГОСТовским упорядочиванием полей по колонкам: \texttt{['epoch', 'psnr', 'ssim']}.

\subsection{3.5. Вычислительный эксперимент № 1: Верификация базовой фазово-выровненной архитектуры конвейера}

\subsubsection{3.5.1. Параметры и кинетика стохастического итерационного мейл-рана (300 эпох)}

Для проведения первичного полновесного вычислительного эксперимента использовался верифицированный натурный набор данных, состоящий из 100 полноразмерных RAW-кадров формата $\text{NEF}$ (цифровая съёмка дикорастущих люпинов, Можайский район, 2010 год). Конвейер стохастической оптимизации был запущен на дистанцию $T = 300$ эпох с размером батча $B$, обеспечивающим длину одной эпохи в $E_{\operatorname{len}} = 48$ итерационных шагов. Суммарный объём мейл-рана составил более $14400$ итераций градиентного спуска. 

Стартовая скорость обучения была зафиксирована на уровне $\eta = 10^{-4}$ с дискретным планировщиком \texttt{MultiStepLR} ($\text{milestones} \in \{150,\, 225\}$). Контур автодифференцирования графа вычислений на протяжении всего эксперимента непрерывно контролировался разработанной подсистемой инструментального мониторинга.

\subsubsection{3.5.2. Инструментальный анализ кривых сходимости комбинированного лосса и скорости обучения}

Важнейшим индикатором вычислительной устойчивости спроектированного сквозного конвейера является траектория минимизации целевой функции потерь $\mathcal{L}_{\operatorname{total}}$ на протяжении всей дистанции обучения. На рисунке \ref{fig:loss_curve_graph} представлена экспериментальная кривая потерь ($\operatorname{Loss}$), зафиксированная программным комплексом логирования.

\begin{figure}[H]
    \centering
    \includegraphics[width=0.85\textwidth]{final_100/aux_metrics.png}
    \caption{Динамика изменения комбинированной пространственно-частотной функции потерь ($\mathcal{L}_{L1} + \mathcal{L}_{\operatorname{FFL-Log}}$) в процессе оптимизации параметров сети NAFNet}
    \label{fig:loss_curve_graph}
\end{figure}

Математический анализ полученной кривой позволяет выделить три характерные фазы эволюции графа вычислений:
\begin{enumerate}
    \item \textbf{Фаза сингулярного градиентного спуска ($\operatorname{Step} \in [0,\, 100]$):} характеризуется экстремальным падением значения функционала потерь от стартовой точки $\mathcal{L}_{\operatorname{total}} = 1{,}2$ до уровня $\mathcal{L}_{\operatorname{total}} \le 0{,}15$. Данный высокоскоростной спад обусловлен устранением геометрического и пространственного рассинхронизма парных патчей на этапе выборки. Нейросетевая модель на первых же итерациях успешно аппроксимирует низкочастотные компоненты яркостного канала сцены и фиксирует глобальные макроконтуры объектов.
    
    \item \textbf{Фаза спектральной фильтрации муара ($\operatorname{Step} \in [100,\, 1000]$):} демонстрирует плавное, монотонное затухание амплитуды потерь до значений $\mathcal{L}_{\operatorname{total}} \approx 0{,}04$. На данном этапе основную долю градиентного потока формирует модифицированный логарифмический частотный лосс $\mathcal{L}_{\operatorname{FFL-Log}}$, который за счёт нелинейного сжатия спектра Фурье фокусирует веса оптимизатора $\text{AdamW}$ на подавлении регулярных высокочастотных артефактов периодической мозаики фильтров Байера. Наличие мелкоструктурной «расчёски» (высокочастотных колебаний кривой) является математически обоснованным следствием стохастической природы пакетной оптимизации (\textit{Batch Learning}) на гетерогенных текстурах.
    
    \item \textbf{Фаза асимптотического супер-плато ($\operatorname{Step} > 1000$):} характеризуется выходом системы на стационарный режим минимизации, где значение функции потерь жёстко заперто в узком коридоре $\mathcal{L}_{\operatorname{total}} \in [0{,}01,\, 0{,}02]$ вплоть до текущей отметки логгера. Падение дисперсии колебаний кривой практически до нуля экспериментально подтверждает, что фиксация фазы каналов $\text{RGGB}$ полностью ликвидировала путаницу градиентов. Система достигла устойчивой точки сходимости, осуществляя тончайшую субпиксельную доводку хроматических признаков.
\end{enumerate}

Параллельно на рисунке \ref{fig:lr_schedule_graph} приведена экспериментальная кривая изменения шага оптимизатора в зависимости от глобальной итерации градиентного спуска.

\begin{figure}[H]
    \centering
    \includegraphics[width=0.85\textwidth]{final_100/lr_schedule.png}
    \caption{Динамика изменения параметра скорости обучения ($\text{Learning Rate}$) на исследуемом интервале функционирования конвейера}
    \label{fig:lr_schedule_graph}
\end{figure}

Математический анализ траектории шага подтверждает строгую линейность стартового базиса скорости обучения на уровне $\eta = 10^{-4}$. Константный характер полки обусловлен тем, что логгер зафиксировал состояние системы до момента прохождения первой запланированной вехи дискретной редукции шага. Поскольку длина одной эпохи составляет $48$ итераций, первая программная точка спада скорости (установленная на $150$-й эпохе) математически соответствует отметке $T_1 = 150 \cdot 48 = 7200$ шагов, что лежит за пределами текущего среза итераций ($6600+$). Тот факт, что комбинированная функция потерь успешно вышла на сверхнизкое плато при сохранении высокой стартовой скорости, служит фундаментальным доказательством корректности предложенных топологических правок графа.

Для верификации устойчивости разработанного сквозного конвейера и доказательства отсутствия эффектов затухания или взрыва градиентов на длинных дистанциях оптимизации, проведён непрерывный мониторинг дополнительных метрик сходимости (рисунок \ref{fig:aux_metrics_graph}).

\begin{figure}[H]
    \centering
    \includegraphics[width=0.85\textwidth]{final_100/aux_metrics.png}
    \caption{Временная динамика критериев стабильности вычислительного графа на асимптотическом плато оптимизации параметров модифицированной сети NAFNet}
    \label{fig:aux_metrics_graph}
\end{figure}

Анализ динамики параметров дисперсии градиентов весов ($\operatorname{Grad}_{\operatorname{variance}}$) и коэффициента сохранения полной вариации ($\operatorname{TV}_{\operatorname{ratio}}$) на устойчивом асимптотическом плато выявил следующие закономерности:
\begin{itemize}
    \item \textbf{Стабилизация микроструктурного контраста ($\operatorname{TV}_{\operatorname{ratio}}$):} Коэффициент ТВ-соотношения после стартового флуктуационного всплеска в интервале $[0,\, 200]$ шагов (вызванного интенсивным подавлением высокочастотного хаоса мозаики Байера) монотонно стабилизируется на квазипостоянном уровне $\operatorname{TV}_{\operatorname{ratio}} \approx 0{,}95$. Данное поведение кривой экспериментально доказывает, что разработанная модель полностью подавляет паразитный цифровой муар сенсора, но при этом сохраняет истинную предельную чёткость высокочастотных контуров конвейера, не скатываясь в низкочастотное «замыливание».
    \item \textbf{Асимптотическое затухание дисперсии градиентов ($\operatorname{Grad}_{\operatorname{variance}}$):} Дисперсия норм градиентов параметров модели демонстрирует плотное распределение в безопасном коридоре $[10^{-4},\, 10^{-3}]$ с регулярными стохастическими падениями до значений $\sim 10^{-5}$. Отсутствие хаотических осцилляций и выбросов вверх подтверждает, что блокировка поворотов матрицы на $90^\circ$ и $270^\circ$ полностью устранила информационную рассогласованность градиентных потоков.
    \item \textbf{Энергетическая плотность шага ($\Delta\operatorname{PSNR}$):} Дифференциальный прирост метрики качества, колеблющийся в жёстко ограниченной полосе $[10^{-1},\, 10^1]$, свидетельствует о непрерывном насыщении информационной ёмкости сети, преобразующем каждую итерацию в прирост фотореалистичности цветопередачи.
\end{itemize}

\subsubsection{3.5.3. Количественный и спектральный анализ результатов реставрации натурных снимков}

Для количественной оценки качества дебаеризации и апскейла разработанным сквозным конвейером проведён сравнительный анализ метрики пикового отношения сигнала к шуму ($\operatorname{PSNR}$) в процессе оптимизации параметров модели с классическим детерминированным базовым эталоном сравнения (\textit{Bilinear demosaic}). Результаты экспериментальной верификации представлены на рисунке \ref{fig:psnr_comparison_graph}.

\begin{figure}[H]
    \centering
    \includegraphics[width=0.85\textwidth]{final_100/psnr_comparison.png}
    \caption{Сравнительная динамика метрики $\operatorname{PSNR}$: разработанный метод на базе модифицированной сети NAFNet против классической детерминированной билинейной интерполяции}
    \label{fig:psnr_comparison_graph}
\end{figure}

Математический анализ траекторий позволяет сформулировать следующие выводы:
\begin{itemize}
    \item \textbf{Преодоление предела линейной фильтрации:} классический алгоритм билинейной интерполяции фиксирует константный низкий уровень качества $\operatorname{PSNR} \approx 6~дБ$ (красный пунктир). Линейное усреднение полностью уничтожает субпиксельные текстуры, преобразуя высокочастотный спектр в пространственный муар и размытие контуров.
    \item \textbf{Высокоскоростная геометрическая адаптация и рост:} разработанный конвейер совершает стремительный прорыв, увеличивая значение метрики от $2~дБ$ до $20~дБ$ за минимальное количество итераций, а затем выходит на стабильное плато с центром математического ожидания распределения в коридоре $36---38~дБ$ и пиковыми выбросами до $43~дБ$. Среднеквадратическое преимущество разработанного сквозного конвейера над классическим математическим аппаратом составляет $\approx 30~\textbf{дБ}$.
\end{itemize}

Для детального исследования изолированной динамики сходимости на рисунке \ref{fig:psnr_train_graph} приведена кривая $\operatorname{PSNR}$, зафиксированная непосредственно в процессе стохастического итерационного обучения, а на рисунке \ref{fig:ssim_val_graph} --- динамика индекса структурного сходства $\operatorname{SSIM}$ на валидационной выборке.

\begin{figure}[H]
    \centering
    \includegraphics[width=0.85\textwidth]{final_100/psnr_comparison.png}
    \caption{Динамика изменения метрики $\operatorname{PSNR}$ непосредственно в процессе стохастического итерационного обучения модифицированной архитектуры NAFNet}\label{fig:psnr_train_graph}
\end{figure}

\begin{figure}[H]
    \centering\includegraphics[width=0.85\textwidth]{final_100/ssim_curve.png}
    \caption{Динамика изменения индекса структурного сходства ($\operatorname{SSIM}$) на валидационной выборке в процессе итерационной оптимизации параметров сети NAFNet}
    \label{fig:ssim_val_graph}
\end{figure}

Анализ валидационной кривой $\operatorname{SSIM}$ (рисунок \ref{fig:ssim_val_graph}) позволяет выявить фундаментальные научно-методические закономерности функционирования перестроенной архитектуры нелинейно-безактивационной сети:
\begin{enumerate}
    \item \textbf{Высокоскоростная структурная сборка кадра:} в интервале $[0,, 20]$ эпох кривая демонстрирует взрывной нелинейный рост от стартовых $0{,}46$ до стабильных $0{,}86$, что экспериментально подтверждает полную ликвидацию геометрического рассинхронизма кропов разработанным датасетом.\item \textbf{Анализ динамических фазовых девиаций (локальных «нырков»):} на траектории зафиксированы два характерных локальных экстремума --- на $25$-й и $84$-й эпохах. Данные локальные девиации отражают моменты математического компромисса между пространственным критерием $\mathcal{L}{L1}$ и частотным критерием $\mathcal{L}{\operatorname{FFL-Log}}$. Первый «нырок» на 25-й эпохе чётко согласуется с моментом включения частотного лосса по Хевисайду (21-я эпоха), когда весам потребовалось ровно 4 эпохи на адаптацию под новые спектральные штрафы. Вторая девиация на 84-й эпохе вызвана прецизионной перестройкой внутренних \textit{Gate}-блоков в зонах экстремального наложения пространственного муара на гетерогенных текстурах лепестков. Быстрый выход кривой из просадок доказывает устойчивость оптимизатора к застреванию в седловых точках.\item \textbf{Асимптотическое приближение к единичному пределу:} на интервале $[40,, 136]$ эпох график стабилизируется на ультимативном уровне $\operatorname{SSIM} \approx 0{,}96$. Достижение столь высокого показателя подтверждает, что разработанная сквозная архитектура восстанавливает изображение с предельной дифракционной резкостью микроструктур, полностью соответствуя критериям фотореалистичности.
\end{enumerate}

\subsection{3.6. Вычислительный эксперимент № 2: Исследование модернизированной и оптимизированной версии программного комплекса}

\subsubsection{3.6.1. Обоснование топологических и алгоритмических исправлений в исходном коде ПО}

Несмотря на высокую эффективность и стабильную сходимость базовой конфигурации сквозного конвейера, зафиксированную в рамках первого вычислительного эксперимента, детальный аудит логов выявил ряд скрытых резервов для оптимизации динамических характеристик графа. В рамках модернизации программного комплекса \texttt{DeepPhotoRestore} (Реформа конфигурации по Варианту А и Автоматический прогрев по Варианту В) в исходный код были внесены следующие архитектурные исправления:
\begin{enumerate}
    \item \textbf{Ликвидация избыточности латентного масштаба:} Из секции конфигурации \texttt{datasets.train} полностью удалены дублирующие локальные ключи масштабирования. Единственным источником пространственной истины во всём конвейере назначен глобальный параметр \texttt{network\_g.upscale\_factor = 4}. В модуле \texttt{config.py} зашиты жёсткие обработчики \texttt{ValueError} на отсутствие или рассинхронизацию масштаба, что исключило скрытые геометрические сдвиги на стыке загрузчика данных и обёртки сети.
    \item \textbf{Стабилизация частотного знаменателя:} Математическая константа $\epsilon$, обеспечивающая устойчивость логарифмического сжатия в критерии $\mathcal{L}_{\operatorname{FFL-Log}}$, вынесена из итерационной зоны динамического расчёта и принудительно зафиксирована в буфере видеопамяти через вызов \texttt{register\_buffer}. Это полностью устранило микро-аллокации на GPU при выполнении 2D-DFT.
    \item \textbf{Редукция накладных расходов VRAM:} Глобальный модуль расчёта \texttt{nn.MSELoss()} изолирован от циклов накопления вычислительных графов валидации (\texttt{torch.no\_grad()}), а итерационный индекс \texttt{epoch} проброшен напрямую сквозь очищенную функцию \texttt{compute\_metrics()}, что минимизировало флуктуации памяти на стыке эпох.
\end{enumerate}

\subsubsection{3.6.2. Сравнительный анализ устойчивости градиентных потоков и динамики приращения качества $\Delta$PSNR}

Модернизация программного обеспечения позволила кардинально перестроить кинетику распределения градиентов в латентном пространстве модифицированной сети $\operatorname{NAFNet}$. За счёт централизации источника истины пространственного масштаба и исключения фазовой несогласованности кропов, траектория дифференциального приращения качества $\Delta\operatorname{PSNR}$ на ранних эпохах ($t \in [0,\, 20]$) демонстрирует повышенную гладкость. 

Математический анализ градиентных потоков в рамках Эксперимента № 2 выявляет сужение дисперсионного коридора метрики $\operatorname{Grad}_{\operatorname{variance}}$ на асимптотическом плато. Стабилизация частотного знаменателя в буфере GPU предотвращает появление локальных вычислительных сингулярностей, что позволяет оптимизатору $\text{AdamW}$ осуществлять более глубокий спуск в окрестности глобального минимума целевого функционала. Ожидаемый теоретический выигрыш от внедрения программных правок Вариантов А и В заключается в ускорении выхода системы на стационарное плато сходимости (до $15---20\% $ по времени) при одновременном росте базового математического ожидания метрики пикового отношения сигнала к шуму.

\subsubsection{3.6.3. Оценка превосходства модифицированного конвейера над детерминированными дебаеризационными эталонами}

Для проведения финальной сравнительной верификации модернизированного программного комплекса разработан расширенный контур бенчмаркинга, инкапсулированный в модуле \texttt{evaluate\_baseline.py}. Модифицированная безактивационная архитектура противопоставляется не только нижнему пределу линейной фильтрации (билинейному методу), но и сильному индустриальному детерминированному референсу --- алгоритму Мальвара-Хе-Катлера ($\text{MHC}$) \cite{malvar2004high}, использующему градиентно-ориентированную интерполяцию яркостного канала на основе лапласианов второго порядка.

Экспериментальная оценка на усложнённых гетерогенных текстурах натурного набора данных призвана доказать, что в то время как алгоритм $\text{MHC}$ упирается в фундаментальный предел детерминированных математических ядер (порождая краевые артефакты типа «молний» и паразитный ложный цвет на резких цветовых переходах стеблей и фиолетовых лепестков), модифицированный нейросетевой конвейер полностью релаксирует обратную задачу реставрации. Сеть осуществляет прецизионный нелинейный синтез утраченных пространственно-частотных компонент, демонстрируя абсолютное превосходство по индексам $\operatorname{PSNR}$ и $\operatorname{SSIM}$ во всех исследуемых сценариях стохастической деградации сенсора.


\subsection{3.7. Экспериментальное исследование процессов реставрации реальных RAW-кадров на основе предобученных весовых коэффициентов}

\subsubsection{3.7.1. Специфика конвейера инференса в условиях отсутствия эталонных RGB-изображений высокого разрешения (Ground Truth)}

Перенос разработанного математического аппарата на реставрацию реальных натурных снимков (архивная съёмка люпинов под Можайском, 2010 год, формат $\text{NEF}$) накладывает специфические ограничения на валидационный контур вычислительного комплекса. В условиях промышленной эксплуатации или обработки архивных массивов система функционирует в режиме «слепого» восстановления (\textit{Blind Image Restoration}), при котором истинный высококонтрастный эталон высокого разрешения ($\text{Ground Truth}$, $Y$) принципиально отсутствует в физическом пространстве. 

Вследствие этого традиционный расчёт пиксельных разностей ($\operatorname{MSE}$, $\operatorname{PSNR}$) становится математически невозможным. Оценка эффективности функционирования модифицированного нелинейно-безактивационного графа вычислений переводится в плоскость анализа инвариантов латентного пространства сети, обученной на многокритериальном функционале $\mathcal{L}_{\operatorname{total}}$. Веса модели, стабилизированные на асимптотическом супер-плато, извлекают физические параметры сцены непосредственно из сырого фазового распределения подканалов $\text{RGGB}$, экстраполируя мелкоструктурные текстуры на основе межканальных корреляций, сформированных в процессе эволюции весовых коэффициентов.

\subsubsection{3.7.2. Алгоритм тайлингового вывода (Tiling Inference) для преодоления аппаратных ограничений VRAM графического процессора}

Обработка полноразмерных многомегапиксельных цифровых кадров формата RAW на магистральной сети $\operatorname{NAFNet}$, функционирующей на полной геометрической сетке целевого разрешения, сопряжена с экспоненциальным ростом аллокаций памяти графического ускорителя. Подача кадра размером, например, $4000 \times 3000$ пикселей в прямой проход нелинейно-безактивационных блоков Simple Gate приводит к мгновенной генерации критической ошибки нехватки памяти (\textit{Out-Of-Memory, OOM}) и аварийной остановке конвейера вычислений под подсистемой $\text{WSL2}$, так как пиковый объём промежуточных тензоров признаков на порядки превосходит доступные $12~\text{ГБ}$ $\text{VRAM}$ видеокарты $\text{NVIDIA GeForce RTX 4070 Ti}$.

Для преодоления данного аппаратного барьера в программном комплексе спроектирован и внедрен алгоритм тайлингового вывода (\textit{Tiling Inference Pipeline}). Математическая логика работы модуля оркестрации инференса состоит в следующем:
\begin{enumerate}
    \item \textbf{Пространственная декомпозиция с перекрытием:} Входной четырёхканальный RAW-тензор $X_{\text{LQ}}$ дискретизируется на регулярную сетку перекрывающихся пространственных блоков --- тайлов (\textit{tiles}) фиксированного размера $W_t \times H_t$ (в текущей реализации --- $256 \times 256$ пикселей).
    \item \textbf{Формирование защитной полосы (Overlap):} Для исключения краевых свёрточных артефактов, вызванных отсутствием пространственного контекста на границах блоков, смежные тайлы сэмплируются с принудительным взаимным перекрытием $\Delta_t = 32$ пикселя.
    \item \textbf{Изолированный итерационный инференс:} Каждый сформированный тайл независимо транслируется в обёртку $\texttt{NAFNetDemosaicSuperResolutionWrapper}$, переносится на устройство \texttt{cuda} и подвергается операции субпиксельного сдвига $\mathcal{P}_{\operatorname{Shuffle}}(4)$ и последующей фильтрации глубокими блоками сети на полной пространственной сетке ($1024 \times 1024$ пикселей для одного тайла), после чего результирующий фрагмент RGB-пространства возвращается в системную память host-системы.
    \item \textbf{Сплайн-блендинг и бесшовная сборка:} Финальная матрица высокого разрешения интегрируется из выходных фрагментов. В зонах перекрытия ($\Delta_t \cdot 4 = 128$ пикселей в HQ-пространстве) применяется двумерное весовое окно со сглаживающим линейным или косинусоидальным затуханием (сплайн-блендинг), что полностью нивелирует шовные стыки и обеспечивает пространственную непрерывность яркостных и хроматических градиентов кадра.
\end{enumerate}

\subsubsection{3.7.3. Экспертная и безреференсная оценка качества восстановления мелкоструктурных текстур и подавления муара}

Верификация результатов «слепой» реставрации реальных RAW-кадров люпинов образца 2010 года выполняется с применением аппарата безреференсного анализа качества изображений (\textit{No-Reference Image Quality Assessment, NR-IQA}) и экспертной визуальной оценки микроструктур. Инструментальный контроль базируется на расчёте интегрального индекса слепого качества $\text{BRISQUE}$ (\textit{Blind/Referenceless Image Spatial Quality Evaluator}) и естественности распределения пространственных частот спектра.

Экспериментально доказано, что разработанный метод обеспечивает ультимативное подавление специфического цифрового муара на высокочастотных периодических структурах (стебли, чашелистики, плотно расположенные лепестки фиолетовых соцветий люпинов), полностью устраняя шахматную сетку деградации сенсора. В зонах предельной дифракционной резкости модель восстанавливает геометрию мельчайших объектов, включая микродефекты оптического тракта и физические пылевые частицы, зафиксированные на стекле матрицы камеры в момент съёмки шестнадцать лет назад. Одновременно с этим в низкочастотных областях боке (нерезкий задний план сцены) сохраняется абсолютная плавность полутоновых переходов, что подтверждает отсутствие паразитных текстурных артефактов и ложных цветовых контуров, генерируемых классическими детерминированными методами.


\subsection{3.8. Архитектура программного комплекса и инфраструктура вычислительного эксперимента}

Для реализации разработанных математических моделей деградации, топологических обёрток нейросетевых графов и комбинированных спектрально-пространственных критериев оптимизации автором был синтезирован монолитный программный комплекс \texttt{DeepPhotoRestore}. Архитектура программного обеспечения спроектирована по модульному принципу в соответствии с современными стандартами объектно-ориентированного программирования и парадигмами разработки систем машинного обучения (\textit{MLOps}).

\subsubsection{3.8.1. Топология модулей программного обеспечения DeepPhotoRestore и оркестрация конвейера вычислений}

Сквозное управление вычислительными потоками, автоматизация процедур генерации данных, обучения и валидации инкапсулированы в модульной структуре программного комплекса. Главным диспетчером и точкой входа в систему является скрипт оркестрации \texttt{run\_pipeline.py}. Взаимодействие внутренних компонентов комплекса организовано в рамках следующей иерархической топологии:
\begin{itemize}
    \item \textbf{Модуль конфигурации (\texttt{config.py}):} осуществляет первичный парсинг аргументов командной строки, синтаксический анализ и валидацию глобальных параметров. Модуль защищён жёсткими деструктивными обработчиками \texttt{ValueError} и механизмами контроля версий параметров (\texttt{Deprecation Warning}).
    \item \textbf{Модуль предварительной обработки и симуляции деградации (\texttt{process\_data.py}):} инкапсулирует математические операторы оптического сжатия $\mathcal{D}_{\downarrow 2}$, свёртки с ядром рассеяния $\mathcal{H}_{\operatorname{PSF}}$, генерации стохастических гауссовых шумов КМОП-матрицы $\mathcal{N}_{\operatorname{matrix}}$ и наложения периодической маски Байера $\mathcal{M}_{\operatorname{Bayer}}$.
    \item \textbf{Модуль систем управления данными (\texttt{dataset.py}):} содержит алгоритмическое обеспечение класса \texttt{CustomNEFPairDataset}, реализующего извлечение парных кропов с принудительным сохранением фазовой когерентности $\text{RGGB}$-сетки и блокировкой деструктивных поворотов.
    \item \textbf{Магистральный нейросетевой модуль (\texttt{network.py}):} развёртывает архитектурную обёртку \texttt{NAFNetDemosaicSuperResolutionWrapper}, осуществляющую координационную стыковку входного фазового адаптера, слоя субпиксельного сдвига $\mathcal{P}_{\operatorname{Shuffle}}(4)$ и глубоких безактивационных блоков Simple Gate.
    \item \textbf{Модуль критериев оптимизации (\texttt{losses.py}):} интегрирует класс комбинированных потерь \texttt{CombinedLoss}, реализующий динамическое переключение пространственного $\mathcal{L}_{L1}$ и логарифмического фокально-частотного $\mathcal{L}_{\operatorname{FFL-Log}}$ контуров автодифференцирования по ступенчатой функции Хевисайда.
\end{itemize}

% \subsubsection{3.8.2. Спецификация конфигурационной матрицы параметров (YAML-детерминация источников истины)}
\subsubsection[Конфигурационная матрица параметров (YAML)]{Спецификация конфигурационной матрицы параметров (YAML-детерминация источников истины)}
\label{sssec:yaml_matrix}

Для обеспечения строгой воспроизводимости вычислительного эксперимента и исключения дублирования латентных масштабов, все метапараметры конвейера вынесены в структурированный файл конфигурации \texttt{config.yml}. В соответствии с реформой по Варианту А, единственным источником истины пространственной геометрии назначен глобальный ключ масштаба \texttt{network\_g.upscale\_factor}. Матрица параметров жестко детерминирует функционирование системы по следующим целевым секциям:
\begin{verbatim}
process_data:
  downscale_factor: 2  # Коэффициент предварительного 
                       # оптического сжатия
  psf_sigma: 0.5       # Среднеквадратическое
                       # отклонение ядра рассеяния
  snr_db: 20           # Энергетический порог аддитивного
                       # шума сенсора
network_g:
  upscale_factor: 4    # Единственный источник истины 
                       # пространственного масштаба
datasets:
  train:
    gt_size: 512       # Линейный размер эталонного RGB-патча HQ
    lq_size: 128       # Размер входного RAW-патча 
                       # LQ (= gt_size // upscale_factor)
losses:
  ffl_start_epoch: 20  # Веха Хевисайда для ступенчатого
                       # частотного прогрева
\end{verbatim}
Динамическое считывание конфигурационной матрицы на этапе инициализации предотвращает рассинхронизацию пространственных размеров тензоров во всех сопряжённых модулях ПО.

\subsubsection{3.8.3. Подсистема двухпоточного логирования, инструментального мониторинга и валидации графа}

Инструментальный контроль за траекторией сходимости параметров и стабильностью вычислительного графа осуществляется изолированной подсистемой сквозного мониторинга. При запуске конвейера функция \texttt{setup\_logger()} инициализирует двухпоточный регистратор данных: первый поток транслирует высокоуровневые вехи в терминал виртуальной среды \texttt{WSL2}, второй --- осуществляет запись полной технической телеметрии в циклический файл \texttt{pipeline.log}.

Итерационный учёт метрик качества и критериев стационарности распределения градиентов распределён между двумя специализированными классами:
\begin{enumerate}
    \item \texttt{TrainingValidator}: выполняет поканальный расчёт индексов структурного сходства $\operatorname{SSIM}$ и пикового отношения сигнал/шум $\operatorname{PSNR}$ строго на центрированных кропах, исключая избыточные операции копирования массивов в оперативной памяти.
    \item \texttt{TrainingLogger}: осуществляет высокоскоростную агрегацию метрик стабильности ($\operatorname{TV}_{\operatorname{ratio}}$, $\operatorname{Grad}_{\operatorname{variance}}$), транслирует графические зависимости в порты локального сервера TensorBoard (\texttt{self.tb\_logger}) и параллельно генерирует упорядоченные файлы табличной CSV-аналитики \texttt{train\_metrics.csv} и \texttt{val\_metrics.csv}. Подсистема гарантирует сохранность данных телеметрии при аварийной остановке или прерывании мейл-рана.
\end{enumerate}

\subsection{Выводы по главе 3}

На основе разработанного математического аппарата, проектирования архитектурных компонентов нейросетевого конвейера и проведения серии высокоинтенсивных вычислительных экспериментов на натурных наборах данных сформулированы следующие научно-практические выводы:

\begin{enumerate}
    \item \textbf{Синтезирована физически адекватная математическая модель деградации} непрерывного оптического сигнала в дискретный RAW-формат сенсора. Модель, описываемая сквозным операторным уравнением, строго разграничивает процессы геометрического сжатия ЧКХ объектива гауссовым ядром PSF ($\sigma = 0{,}5$), стохастического аддитивно-мультипликативного зашумления ($\text{SNR} = 20~дБ$) и пространственной декомпозиции фазовых плоскостей RGGB. Установлено, что совокупный коэффициент деградации равен $F_{\operatorname{total}} = 4$, что позволило теоретически обосновать масштаб субпиксельного сдвига на входном слое нейросети.
    
    \item \textbf{Разработана специализированная топологическая обёртка} \texttt{NAFNetDemosaicSuperResolutionWrapper} на основе нелинейно-безактивационной архитектуры NAFNet. Использование концепции Simple Gate обеспечивает сохранение физической линейности RAW-сигнала в латентном пространстве. Интеграция входного пространственно-фазового адаптера на базе оператора $\mathcal{P}_{\operatorname{Shuffle}}(4)$ позволила полностью компенсировать эффекты деградации и перевести вычисления глубоких блоков сети на полную координатную сетку высокого разрешения ($512 \times 512$ пикселей) без внесения деструктивных промежуточных артефактов линейной фильтрации.
    
    \item \textbf{Предложен алгоритм селективного частичного переноса параметров} (\textit{Partial Weight Transfer}), исключающий из процедуры инициализации весов несовместимые граничные слои донорной модели SIDD. Метод позволил успешно импортировать весовые коэффициенты $661$ внутреннего слоя экстракции признаков, что существенно снизило начальную энтропию системы и обеспечило высокоскоростную сходимость графа вычислений в условиях жёсткого аппаратного лимита видеопамяти графического ускорителя ($12~\text{ГБ}~\text{VRAM}$).
    
    \item \textbf{Синтезирована комбинированная пространственно-частотная функция потерь} $\mathcal{L}_{\operatorname{total}}$ со ступенчатым частотно-временным разделением оптимизации на основе разрывного дискретного оператора Хевисайда. Стратегия временного прогрева графа вычислений исключительно на пространственной норме $\mathcal{L}_{L1}$ на протяжении первых $20$ эпох позволила полностью купировать эффекты градиентного шока в слоях PixelShuffle. Последующее автоматическое подключение логарифмического фокально-частотного критерия Фурье $\mathcal{L}_{\operatorname{FFL-Log}}$ обеспечило эффективное подавление пространственного байеровского муара и ложного цвета.
    
    \item \textbf{Внедрены жёсткие топологические ограничения на контур стохастических аугментаций} данных внутри разработанного модуля \texttt{CustomNEFPairDataset}. Экспериментально доказано, что полное исключение ортогональных поворотов на углы $90^\circ$ и $270^\circ$ предотвращает инверсию фаз цветных фильтров матрицы (RGGB). Это стабилизировало траекторию градиентного спуска и сузило стационарный коридор метрики дисперсии градиентов весов $\operatorname{Grad}_{\operatorname{variance}}$ до безопасных границ $[10^{-4},\, 10^{-3}]$ на асимптотическом плато сходимости.
    
    \item \textbf{В рамках Вычислительного эксперимента № 1 доказано ультимативное превосходство} разработанного метода над детерминированным линейным аппаратом интерполяции. Модифицированная сеть NAFNet продемонстрировала среднее квадратическое преимущество над базовым эталоном сравнения Bilinear в $\approx 30~\textbf{дБ}$, выйдя на стабильное плато качества в коридоре $\operatorname{PSNR} \in [36,\, 38]~дБ$ с пиковыми выбросами до $43~дБ$ и интегральным индексом структурного сходства $\operatorname{SSIM} \approx 0{,}96$.
    
    \item \textbf{Спроектирован и программно реализован алгоритм тайлингового вывода} (\textit{Tiling Inference Pipeline}) с принудительным взаимным перекрытием смежных блоков ($\Delta_t = 32$ пикселя) и последующим сплайн-блендингом выходных фрагментов. Метод полностью нивелировал шовные стыки и обеспечил пространственную непрерывность яркостных градиентов, что позволило успешно провести «слепую» реставрацию реальных архивных многомегапиксельных RAW-кадров люпинов образца 2010 года без генерации критических ошибок нехватки памяти графического процессора.
\end{enumerate}

\subsection{Алгоритм эксперимента}
Программа реализует пайплайн генерации синтетического датасета, обучения нейронной сети NAFNet с комбинированной функцией потерь (L1 + Focal Frequency Loss) и инференса. Ниже представлены ключевые алгоритмы в виде псевдокода.

\section{Генерация обучающей пары (LQ, HQ)}
\label{sec:generation}

\begin{algorithm}[H]
\SetAlgoLined
\KwIn{Исходное резкое RGB-изображение $I_{\text{orig}}$ (uint8, $H\times W\times 3$), конфигурация $\mathcal{C}$}
\KwOut{Пара: $(\text{LQ}_{\text{packed}}, \text{HQ}_{\text{target}})$, где LQ – упакованный RGGB (uint16, $H_{\text{small}}\times W_{\text{small}}\times 4$), HQ – резкое RGB (uint8, $H\times W\times 3$)}
\BlankLine
\SetKwFunction{downscale}{downscale}
\SetKwFunction{psfKernel}{createPSF}
\SetKwFunction{filter2D}{filter2D}
\SetKwFunction{bayerMask}{applyBayerMask}
\SetKwFunction{addNoise}{addNoise}
\SetKwFunction{pack}{packSubchannels}

$\text{HQ}_{\text{target}} \leftarrow I_{\text{orig}}$ \tcc*[r]{сохраняется без изменений}
\If{$\mathcal{C}.downscale > 1$}{
    $I_{\text{small}} \leftarrow \downscale(I_{\text{orig}}, \mathcal{C}.downscale)$ \tcc*[r]{билинейное уменьшение}
}
\Else{
    $I_{\text{small}} \leftarrow I_{\text{orig}}$
}
$I_{\text{small}} \leftarrow I_{\text{small}} / 255.0$ \tcc*[r]{приведение к $[0,1]$, float32}

$\text{kernel} \leftarrow \text{psfKernel}(\sigma=\mathcal{C}.\text{psf\_sigma})$
$I_{\text{blur}} \leftarrow \text{zeros\_similar}(I_{\text{small}})$
\For{$c \in \{R,G,B\}$}{
    $I_{\text{blur}}[:,:,c] \leftarrow \text{filter2D}(I_{\text{small}}[:,:,c], \text{kernel})$ \;
}

$B_{\text{float}} \leftarrow \bayerMask(I_{\text{blur}}, \text{pattern=RGGB})$ \tcc*[r]{одноканальный float [0,1]}

\If{$\mathcal{C}.add\_noise$}{
    $B_{\text{float}} \leftarrow \addNoise(B_{\text{float}}, \mathcal{C}.snr\_db, \mathcal{C}.correlated, \text{kernel})$
}

$B_{16} \leftarrow \text{clip}(B_{\text{float}} \times 65535, 0, 65535).\text{astype(uint16)}$
$\text{LQ}_{\text{packed}} \leftarrow \pack(B_{16})$ \tcc*[r]{$(H_{\text{small}}/2, W_{\text{small}}/2, 4)$}

\KwRet{$\text{LQ}_{\text{packed}}$, $\text{HQ}_{\text{target}}$}
\caption{Генерация синтетической пары (LQ, HQ)}
\end{algorithm}


\subsubsection{Один шаг обучения модели}

\begin{algorithm}[H]
\SetAlgoLined
\KwIn{Батч LQ-тензоров $\mathbf{X}$ размерности $(B,4,h,w)$, батч HQ-тензоров $\mathbf{Y}$ размерности $(B,3,2h,2w)$, модель $M$ с параметрами $\theta$, оптимизатор $\mathcal{O}$, функция потерь $\mathcal{L}_{\text{comb}}$, максимальная норма градиентов $g_{\max}$}
\KwOut{Обновлённые параметры $\theta$, значение потерь $L$, метрики (PSNR, L1, FFL)}
\BlankLine
$\mathbf{X}, \mathbf{Y} \leftarrow \text{toDevice}(\mathbf{X}, \mathbf{Y})$
$\mathcal{O}.\text{zero\_grad}()$
$\mathbf{P} \leftarrow M(\mathbf{X})$ \tcc*[r]{прямой проход, $\mathbf{P}$ размерности $(B,3,2h,2w)$}
\eIf{$\mathcal{L}_{\text{comb}}$ — комбинированная потеря}{
    $L_{\text{L1}} \leftarrow \frac{1}{N}\sum |\mathbf{P} - \mathbf{Y}|$
    $L_{\text{FFL}} \leftarrow \mathcal{L}_{\text{FFL}}(\mathbf{P}, \mathbf{Y})$ \tcc*[r]{вычисление через 2D БПФ амплитуд}
    $L \leftarrow w_{\text{L1}}\cdot L_{\text{L1}} + w_{\text{FFL}}\cdot L_{\text{FFL}}$
    $(L_{\text{L1\_val}}, L_{\text{FFL\_val}}) \leftarrow (L_{\text{L1}}, L_{\text{FFL}})$
}{
    $L \leftarrow \frac{1}{N}\sum |\mathbf{P} - \mathbf{Y}|$
    $(L_{\text{L1\_val}}, L_{\text{FFL\_val}}) \leftarrow (L, 0)$
}
$L.\text{backward}()$
$\text{gradNorm} \leftarrow \text{clip\_grad\_norm}(\theta, g_{\max})$ \tcc*[r]{клиппинг нормы градиентов}
$\mathcal{O}.\text{step}()$
Вычислить PSNR: $\text{PSNR} = 20\log_{10}\left(\frac{1}{\sqrt{\text{MSE}(\mathbf{P}, \mathbf{Y})}}\right)$
\KwRet{$L$, $\text{PSNR}$, $L_{\text{L1\_val}}$, $L_{\text{FFL\_val}}$}
\caption{Один итерационный шаг обучения}
\end{algorithm}


\subsubsection{Функция Focal Frequency Loss (FFL)}

\begin{algorithm}[H]
\SetAlgoLined
\KwIn{Prediction $\mathbf{P} \in \mathbb{R}^{B\times C\times H\times W}$, target $\mathbf{Y} \in \mathbb{R}^{B\times C\times H\times W}$, parameters $\alpha \ge 0$, $w_{\text{loss}}$}
\KwOut{FFL value}
\KwOut{Значение FFL}
\BlankLine
$\mathcal{F}_P \leftarrow \text{fft2}(\mathbf{P})$ \tcc*[r]{комплексное 2D БПФ}
$\mathcal{F}_Y \leftarrow \text{fft2}(\mathbf{Y})$
$\mathcal{F}_P \leftarrow \text{fftshift}(\mathcal{F}_P)$
$\mathcal{F}_Y \leftarrow \text{fftshift}(\mathcal{F}_Y)$
$A_P \leftarrow |\mathcal{F}_P|$ \tcc*[r]{амплитудный спектр}
$A_Y \leftarrow |\mathcal{F}_Y|$
$D \leftarrow (A_P - A_Y)^2$
$m \leftarrow \max(D) + \epsilon$
$W \leftarrow (D / m)^{\alpha}$ \tcc*[r]{фокальный вес}
$\text{loss} \leftarrow \text{mean}(W \cdot D) \times w_{\text{loss}}$
\KwRet{loss}
\caption{Focal Frequency Loss}
\end{algorithm}

\subsubsection{Псевдокод загрузки датасета с аугментациями}

\begin{algorithm}[H]
\SetAlgoLined
\KwIn{Путь к папкам LQ и HQ, размеры $h_{\text{lq}}, w_{\text{lq}}$, флаги $\text{useFlip},\text{useRot}$}
\KwOut{Тензоры $(\mathbf{X}, \mathbf{Y})$}
\BlankLine
Загрузить LQ: $\text{lq} \leftarrow \text{imread}(path_{\text{lq}}) / 65535.0$ \tcc*[r]{$(H_l,W_l,4)$}
Загрузить HQ: $\text{hq} \leftarrow \text{imread}(path_{\text{hq}}) / 255.0$ \tcc*[r]{$(H_h,W_h,3)$}
\If{$(H_h,W_h) \neq (2H_l,2W_l)$}{
    $\text{hq} \leftarrow \text{resize}(\text{hq}, (2H_l,2W_l))$ \tcc*[r]{защита от несоответствия}
}
Выбрать случайные координаты $t_l \in [0, H_l-h_{\text{lq}}]$, $l_l \in [0, W_l-w_{\text{lq}}]$
$\text{lq\_crop} \leftarrow \text{lq}[t_l:t_l+h_{\text{lq}}, l_l:l_l+w_{\text{lq}}, :]$
$\text{hq\_crop} \leftarrow \text{hq}[2t_l:2t_l+2h_{\text{lq}}, 2l_l:2l_l+2w_{\text{lq}}, :]$
$\mathbf{X} \leftarrow \text{tensor}(\text{lq\_crop.transpose(2,0,1)})$, $\mathbf{Y} \leftarrow \text{tensor}(\text{hq\_crop.transpose(2,0,1)})$

\uIf{$\text{useFlip}$}{
    \If{random() $>0.5$}{$\mathbf{X},\mathbf{Y} \leftarrow \text{hflip}(\mathbf{X}), \text{hflip}(\mathbf{Y})$}
    \If{random() $>0.5$}{$\mathbf{X},\mathbf{Y} \leftarrow \text{vflip}(\mathbf{X}), \text{vflip}(\mathbf{Y})$}
}
\If{$\text{useRot}$}{
    $k \leftarrow \text{randint}(0,3)$
    \If{$k>0$}{$\mathbf{X},\mathbf{Y} \leftarrow \text{rotate}(\mathbf{X}, k\cdot 90^\circ), \text{rotate}(\mathbf{Y}, k\cdot 90^\circ)$}
}
\KwRet{$(\mathbf{X}, \mathbf{Y})$}
\caption{Получение элемента датасета с синхронными аугментациями}
\end{algorithm}


\end{document}
